\begin{definition}
\label{definition-regular-sequence}
Let $R$ be a ring. Let $M$ be an $R$-module. A sequence of elements
$f_1, \ldots, f_r$ of $R$ is called an {\it $M$-regular sequence}
if the following conditions hold:
\begin{enumerate}
\item $f_i$ is a nonzerodivisor on
$M/(f_1, \ldots, f_{i - 1})M$
for each $i = 1, \ldots, r$, and
\item the module $M/(f_1, \ldots, f_r)M$ is not zero.
\end{enumerate}
If $I$ is an ideal of $R$ and $f_1, \ldots, f_r \in I$
then we call $f_1, \ldots, f_r$ an {\it $M$-regular sequence
in $I$}. If $M = R$, we call $f_1, \ldots, f_r$ simply a
{\it regular sequence} (in $I$).
\end{definition}